\section{Evaluating Language Models: Perplexity}

\begin{frame}[allowframebreaks]{Perplexity}

\begin{itemize}
    \item Perplexity (PPL) measures how well a language model predicts a sequence of words.
    Intuitively, it reflects ``on average, how many choices the model is unsure about'' per word.

    \item Formula:
    \[
        \text{PPL}(W) = P(w_1, \dots, w_N)^{-\frac{1}{N}}
        = \exp\Bigg(-\frac{1}{N} \sum_{i=1}^{N} \log P(w_i \mid w_1, \dots, w_{i-1}) \Bigg)
    \]

    \item How to interpret:
    \begin{itemize}
        \item Lower perplexity → model assigns higher probability to the test set sequences → better predictions.
    \end{itemize}
\framebreak
    \item Why it matters:
    \begin{itemize}
        \item Standard metric for evaluating language models, from n-gram models to neural LMs.
        \item Only compare perplexities of models that share the same vocabulary $V$.
    \end{itemize}

    \item Example, scoring both models on the same held-out text with the same $V$:
    \begin{itemize}
        \item Bigram model → PPL $\approx$ 200: on average it is choosing between
        roughly two hundred plausible next words.
        \item Neural LM → PPL $\approx$ 20: the same choice, narrowed tenfold.
    \end{itemize}

    \item Caveat: perplexity rewards fluency on held-out text, and nothing else. It
    says nothing about whether the model is truthful, useful, or safe.
\end{itemize}

\end{frame}
